% from ega0-23-fr.tex
\oldpage[0\textsubscript{IV-1}]{217}

\emph{(ii)} L'hypothèse entraîne que les $L_i$ sont des extensions séparables
de $K$, donc $K'\otimes_KL_i$ est réduit (Bourbaki, \emph{Alg.}, chap.~VIII,
§~7, n\textsuperscript{o}~3, th.~1) et l'on peut appliquer (i) à toute
extension finie $K'$ de $K$, ce qui démontre notre assertion.
\end{proof}

\subsection{Clôture intégrale d'un anneau local noethérien intègre}
\label{subsection:0.23.2-fr}

\begin{env}[23.2.1]
\label{0.23.2.1-fr}
Dans ce qui suit, pour tout anneau $A$, nous noterons $A_{\mathrm{red}}$ pour
simplifier, le quotient de $A$ par son nilradical (de sorte que si
$X=\operatorname{Spec}(A)$, on a
$X_{\mathrm{red}}=\operatorname{Spec}(A_{\mathrm{red}})$).

Nous dirons qu'un anneau \emph{local} $A$ est \emph{unibranche} si l'anneau
$A_{\mathrm{red}}$ est intègre et si la clôture intégrale de
$A_{\mathrm{red}}$ est un anneau \emph{local}, ce qui généralise la définition
donnée dans (III, 4.3.6). Nous dirons que $A$ est \emph{géométriquement
unibranche} s'il est unibranche et si le corps résiduel de l'anneau local,
clôture intégrale de $A_{\mathrm{red}}$, est une extension \emph{radicielle}
de celui de $A$. Il est clair qu'un anneau local intégralement clos est
géométriquement unibranche~; il résulte de (23.1.6) qu'un anneau local
noethérien intègre \emph{complet} est unibranche.
\end{env}

\begin{lemma}[23.2.2]
\label{0.23.2.2-fr}
\textnormal{(*)}\footnote{Dans la suite de ce numéro, nous utilisons des
notions et résultats exposés au chap.~IV, §§~2 et 5~; comme les résultats de
ce numéro ne sont pas utilisés avant le chap.~IV, §~6, cela n'entraîne pas de
cercle vicieux.} Soient $A$ un anneau intègre, $A'$ sa clôture intégrale~; pour
que le morphisme canonique
$f:\operatorname{Spec}(A')\to\operatorname{Spec}(A)$ soit radiciel, il faut et
il suffit que pour tout $\mathfrak p\in\operatorname{Spec}(A)$,
$A_{\mathfrak p}$ soit géométriquement unibranche~; lorsqu'il en est ainsi,
$f$ est un homéomorphisme.
\end{lemma}

\begin{proof}
En effet, pour tout $\mathfrak p\in\operatorname{Spec}(A)$, la clôture
intégrale de $A_{\mathfrak p}$ est $A'_{\mathfrak p}$ (Bourbaki,
\emph{Alg. comm.}, chap.~V, §~1, n\textsuperscript{o}~5, prop.~16), et tous
les idéaux premiers de $A'_{\mathfrak p}$ au-dessus de
$\mathfrak pA_{\mathfrak p}$ sont maximaux (\emph{loc. cit.}, §~2,
n\textsuperscript{o}~1, prop.~1). Dire que $f$ est injectif signifie donc que
pour tout $\mathfrak p\in\operatorname{Spec}(A)$, $A'_{\mathfrak p}$ est un
anneau \emph{local}, c'est-à-dire que les $A_{\mathfrak p}$ sont
unibranches. Dire que tout $A_{\mathfrak p}$ est géométriquement unibranche
signifie alors que $f$ est radiciel en vertu de (I, 3.5.8). Lorsqu'il en est
ainsi, $f$ est surjectif et fermé (II, 6.1.10), donc un homéomorphisme.
\end{proof}

\begin{lemma}[23.2.3]
\label{0.23.2.3-fr}
Soient $A$ un anneau intègre, $K$ son corps des fractions, $B$ un anneau
noethérien, $\varphi:A\to B$ un homomorphisme d'anneau faisant de $B$ un
$A$-module plat. Soit $K'$ une extension de $K$~; posons
$B_1=B_{\mathrm{red}}$, et soit $R$ l'anneau total des fractions de $B_1$.
Alors, pour toute sous-$A$-algèbre $A'$ de $K'$, l'homomorphisme canonique
\begin{equation}
\label{0.23.1.8.1-fr}
  (A'\otimes_AB)_{\mathrm{red}}
  \longrightarrow(K'\otimes_AR)_{\mathrm{red}}
  \tag{23.1.8.1}
\end{equation}
est injectif.
\end{lemma}

\begin{proof}
On peut en effet considérer l'homomorphisme canonique
$h:A'\otimes_AB\to K'\otimes_AR$ comme composé des homomorphismes suivants
\[
  A'\otimes_AB\xrightarrow{u}K'\otimes_AB
  \xrightarrow{v}K'\otimes_AB_1\xrightarrow{w}K'\otimes_AR.
\]
Comme $B$ est un $A$-module plat, $u$ est injectif~; de même, $K'$ étant un
$K$-module plat et $K$ un $A$-module plat, $K'$ est un $A$-module plat, et
$w$ est donc injectif puisque $B_1$ est réduit, donc l'homomorphisme
$B_1\to R$ injectif. Enfin, pour la même raison, si $\mathfrak N$ est

